\documentclass[11pt]{article}

\usepackage[margin=1in]{geometry}
\usepackage{amsmath}
\usepackage{amssymb}
\usepackage{graphicx}
\usepackage{booktabs}
\usepackage{multirow}
\usepackage{array}
\usepackage{hyperref}
\usepackage{url}
\usepackage{natbib}
\usepackage{xcolor}
\usepackage{listings}
\usepackage{caption}
\usepackage{subcaption}
\usepackage{float}
\usepackage{enumitem}
\usepackage{titlesec}
\usepackage{microtype}

\hypersetup{
    colorlinks=true,
    linkcolor=blue,
    citecolor=blue,
    urlcolor=blue
}

\lstset{
    basicstyle=\ttfamily\small,
    breaklines=true,
    frame=single,
    backgroundcolor=\color{gray!10}
}

\title{\textbf{Automated Scoring of Pharmacy Student Reflections:\\
A Decomposed LLM Pipeline with Cross-Model Consistency Validation}}

\author{
    Bella Chang \and Emily Wang \and Kristina Fujimoto\\[0.5em]
    %\textit{New York University and University of Washington}
}

\date{}

\begin{document}

\maketitle

\begin{abstract}
Reflective writing is a core component of pharmacy education, helping students synthesize clinical experiences and develop professional identity. However, manually scoring reflections at scale is time-consuming and subject to inter-rater variability. We present an automated scoring pipeline for pharmacy student reflections structured around a three-dimensional rubric derived from course requirements at the University of Washington School of Pharmacy: \textbf{WHAT} (description of learning), \textbf{WHY} (analysis of significance), and \textbf{HOW} (integration and actionability). We investigate multiple computational approaches—TF-IDF with logistic regression, sentence embeddings with logistic regression, and large language model (LLM) prompting across the Claude, GPT-4o, and Llama families—comparing decomposed per-dimension scoring against direct total score prediction. Based on experimental evaluation, we adopt an LLM-based final pipeline using Claude Sonnet as the primary grader across all three dimensions, with Llama~3.3-70B serving as a secondary cross-model validator. Evidence extraction grounds each prediction in exact text spans from the student's reflection, and reflections with low cross-model agreement are flagged for instructor review. We discuss limitations arising from data constraints and model cost, and outline directions for future work including broader data collection and expanded model coverage.
\end{abstract}

% ─────────────────────────────────────────────
\section{Introduction}
% ─────────────────────────────────────────────

Reflective practice is central to health professions education. In pharmacy programs, students are regularly asked to write structured reflections on clinical encounters, coursework, and professional development experiences. These reflections are assessed by instructors according to rubrics that probe the depth of a student's description, analysis, and forward-looking integration of their learning. While rubric-based scoring provides a shared evaluative framework, manually reviewing dozens or hundreds of reflections places a significant burden on faculty, particularly in large cohorts.

Automated scoring of short-answer and reflective text has been an active area of NLP research \citep{burstein2003automated, shermis2013handbook}. Recent advances in large language models (LLMs) have made it possible to prompt a model with a rubric and receive structured, reasoned evaluations \citep{brown2020language, wei2022chain}. At the same time, classical approaches such as TF-IDF bag-of-words representations and sentence embedding classifiers offer computationally cheaper alternatives for benchmarking, allowing us to assess whether the added cost and complexity of LLMs is warranted for this task \citep{reimers2019sentence}.

This paper describes the design, experimentation, and final architecture of an automated reflection scoring system developed in collaboration with faculty at the University of Washington School of Pharmacy. Our contributions are:

\begin{itemize}[leftmargin=1.5em]
    \item A systematic comparison of TF-IDF, sentence embedding, and LLM-based scoring across three rubric dimensions using both decomposed and direct scoring strategies.
    \item An LLM-based final pipeline using Claude Sonnet across all three rubric dimensions, selected based on empirical performance and practical cost considerations.
    \item An evidence extraction and verification module that grounds each score in exact text spans from the student's reflection.
    \item A cross-model consistency validation protocol using Llama~3.3-70B for auditing low-agreement cases.
\end{itemize}

% ─────────────────────────────────────────────
\section{Related Work}
% ─────────────────────────────────────────────

\paragraph{Automated Essay and Short-Answer Scoring.}
Automated scoring of student writing has a long history, from early systems based on lexical features \citep{page1966imminence} to modern neural approaches \citep{taghipour2016neural}. Most work targets holistic essay scoring; rubric-aligned dimension scoring is less common but has been explored in educational contexts \citep{mathias2020happy}.

\paragraph{LLMs for Educational Assessment.}
GPT-family and instruction-tuned models have been applied to grade short answers, generate feedback, and align scores to rubrics \citep{kortemeyer2023could, mizumoto2023exploring}. Chain-of-thought prompting improves scoring reliability by eliciting explicit reasoning before numeric output \citep{wei2022chain}.

\paragraph{Sentence Embeddings for Classification.}
Sentence transformer models such as \texttt{all-MiniLM-L6-v2} \citep{reimers2019sentence} produce dense semantic representations well-suited for downstream classification. We include embedding-based classifiers in our experimental comparison as a strong, cost-efficient baseline against which LLM performance can be benchmarked.

\paragraph{Reflective Writing Assessment.}
Automated analysis of reflective writing has been explored in nursing \citep{song2022automated} and medical education \citep{kennedy2014reflection}. These works highlight that reflection quality is multidimensional and that surface-level lexical cues are insufficient proxies for analytic depth—motivating our adoption of LLM reasoning as the core of the final pipeline.

% ─────────────────────────────────────────────
\section{Task Formulation}
% ─────────────────────────────────────────────

\subsection{Rubric Dimensions}

The scoring rubric was derived directly from the reflection assessment requirements of a pharmacy course at the University of Washington School of Pharmacy, in collaboration with course faculty. Each student reflection is scored on three dimensions, each on an ordinal scale of 0--2 (see Appendix~\ref{app:rubric} for the full rubric as provided by faculty):

\begin{description}[leftmargin=2em, labelindent=0em]
    \item[\textbf{WHAT (Description).}] Does the student clearly describe a specific learning experience, skill, or concept?
    \begin{itemize}
        \item \textbf{2}: Clear, specific description with concrete and relevant details.
        \item \textbf{1}: Vague, incomplete, or lacking specificity.
        \item \textbf{0}: No meaningful description, or off-topic.
    \end{itemize}

    \item[\textbf{WHY (Analysis).}] Does the student explain why this learning is personally significant?
    \begin{itemize}
        \item \textbf{2}: Clear personal significance with insight, self-awareness, and concrete examples.
        \item \textbf{1}: Surface-level or generic analysis, lacking personal connection.
        \item \textbf{0}: No meaningful analysis of importance.
    \end{itemize}

    \item[\textbf{HOW (Integration).}] Does the student articulate a specific, actionable plan for applying the learning?
    \begin{itemize}
        \item \textbf{2}: Specific, actionable, and realistic plan for future pharmacy practice.
        \item \textbf{1}: Vague, generic, or not clearly actionable.
        \item \textbf{0}: No mention of future application, or purely hypothetical.
    \end{itemize}
\end{description}

The total score is the sum of the three dimension scores, ranging from 0 to 6.

% ─────────────────────────────────────────────
\section{Dataset}
% ─────────────────────────────────────────────

\subsection{Data Collection and Privacy Constraints}

We worked in collaboration with University of Washington School of Pharmacy faculty to design a rubric and scoring protocol. Due to student data privacy regulations and IRB constraints, we were unable to directly use real student reflections. Our work proceeded in two phases:

\begin{enumerate}
    \item \textbf{Initial development dataset.} We first created a small self-authored dummy dataset to verify the preprocessing and scoring pipeline infrastructure. This dataset was not validated by faculty and served only as a functional check.

    \item \textbf{Expanded synthetic dataset.} In collaboration with faculty, we constructed an expanded synthetic dataset of reflections representing the full range of score combinations across WHAT, WHY, and HOW. Faculty reviewed and validated this dataset to confirm that the synthetic reflections reflected realistic student language patterns at each score level. This dataset served as the basis for Experiment~2 (cross-model evaluation) and as the final runthrough set for the deployed pipeline. A derived subset of this dataset was additionally used for Experiment~1 (decomposed vs.\ direct scoring).
\end{enumerate}

Each reflection in the synthetic datasets is annotated with ground-truth WHAT, WHY, and HOW scores. Both datasets are small (a practical limitation discussed in Section~\ref{sec:limitations}), but the expanded dataset is structurally representative of the rubric's score space.

\subsection{Exploratory Data Analysis}

Before modeling, we conducted exploratory data analysis (EDA) to characterize the dataset along several axes.

\paragraph{Score Distributions.} We examine the frequency of each score level (0, 1, 2) per dimension to detect class imbalance, which motivates the use of balanced class weighting in classical classifiers and informs how we interpret agreement metrics across models.

\paragraph{Length Distributions.} Word count and sentence count are examined by total score quartile. Higher-quality reflections tend to be longer, though length alone is a weak predictor.

\paragraph{Language Marker Analysis.} We define four families of lexical markers hypothesized to correlate with rubric dimensions:
\begin{itemize}
    \item \textit{Specificity markers} (e.g., ``specifically'', ``patient'', ``medication'', ``clinical'') — associated with WHAT.
    \item \textit{Causal markers} (e.g., ``because'', ``therefore'', ``as a result'') — associated with WHY.
    \item \textit{Future markers} (e.g., ``will'', ``plan'', ``intend'', ``commit'') — associated with HOW.
    \item \textit{Introspective markers} (e.g., ``I realize'', ``I learned'', ``I believe'') — associated with WHY depth.
\end{itemize}
Marker rate (count normalized by word count) is averaged per score level to validate these hypotheses. This analysis is useful for instructors as a diagnostic tool: if markers associated with a given dimension do not vary meaningfully across score levels, it may indicate that the rubric language is not well-captured by surface lexical signals—informing decisions about how to structure prompts or rubric descriptors.

\paragraph{Sentiment Analysis.} We apply VADER \citep{hutto2014vader} to compute compound sentiment scores per reflection and examine their distribution by total score. The intuition behind this analysis is to understand the affective character of student reflections: are higher-quality reflections more positive in tone, perhaps reflecting greater engagement or optimism about future practice? Conversely, low-quality reflections might be more neutral or perfunctory. This can help instructors understand whether sentiment correlates with depth of engagement, which may be a useful secondary signal when reviewing borderline cases.

% ─────────────────────────────────────────────
\section{Scoring Approaches}
% ─────────────────────────────────────────────

We implement and compare three families of scoring approaches, beginning with classical baselines before moving to LLM-based methods.

\subsection{Approach A: TF-IDF with Logistic Regression}

As a classical baseline, we represent each reflection as a TF-IDF vector with up to 5{,}000 features, using unigrams and bigrams with English stopword removal. A logistic regression classifier with balanced class weighting is trained per target dimension using stratified 5-fold cross-validation. We include this approach to establish a lower bound on performance and to assess how much of the scoring signal can be captured by surface-level lexical patterns alone. This serves as a useful reference point: if TF-IDF achieves comparable performance to LLMs, it would suggest that the rubric dimensions are largely distinguishable by vocabulary alone. As we show in our experiments, this is not the case for all dimensions—motivating the use of richer representations.

\subsection{Approach B: Sentence Embeddings with Logistic Regression}

We encode each reflection using \texttt{all-MiniLM-L6-v2} \citep{reimers2019sentence}, a compact sentence transformer that produces 384-dimensional dense embeddings capturing semantic meaning beyond surface vocabulary. A logistic regression head is trained per dimension on these embeddings with the same cross-validation protocol as Approach A. We include this approach as an intermediate baseline: it is more semantically expressive than TF-IDF and substantially cheaper than LLM inference, making it a practical alternative if LLM performance is not sufficiently better to justify the cost.

\subsection{Approach C: LLM Prompting}

We prompt large language models with the full rubric specification and request structured JSON output containing a score, an exact evidence span from the reflection, and a textual explanation for each dimension. The prompt is designed to elicit chain-of-thought style reasoning before numeric output. We evaluate models from three families, chosen to cover a range of sizes, providers, and access modalities:

\begin{itemize}
    \item \textbf{Claude family} (Anthropic): Claude Haiku and Claude Sonnet. We test both to assess whether larger-within-family models improve scoring. Claude is a strong candidate given its strong instruction following, structured JSON output, and documented reasoning capabilities.
    \item \textbf{GPT family} (OpenAI): GPT-4o. Included as the dominant alternative to Claude in the academic and industry community, enabling cross-provider comparison.
    \item \textbf{Llama family} (Meta, via Groq API): Llama 3.3-70B and Llama 4 Scout. We include Llama models as a representative open-weight family. Open-weight models are particularly relevant for future privacy-preserving deployment scenarios (local inference without transmitting student data to third-party APIs), and the Groq API provided a cost-effective way to evaluate them at scale.
\end{itemize}

LLMs are queried at temperature 0 for all scoring to maximize reproducibility. Automatic retry logic handles JSON parsing failures.

% ─────────────────────────────────────────────
\section{Experiments}
% ─────────────────────────────────────────────

\noindent\textit{Note on scoring strategy:} A key design question is whether to predict each dimension independently and sum (decomposed scoring) or to predict the total score directly (direct scoring). Decomposed scoring allows different models to be applied per dimension and produces interpretable per-dimension feedback. We treat this as an empirical question in Experiment~1 before evaluating individual approaches in Experiment~2.

\subsection{Experiment 1: Decomposed vs.\ Direct Scoring}

\paragraph{Motivation.} Decomposed scoring requires training three separate classifiers (one per dimension) instead of one, adding complexity. We test whether this decomposition is justified by measuring whether it improves total score prediction—particularly on QWK, which penalizes large ordinal errors.

\paragraph{Setup.} Using Approaches A and B (TF-IDF and embeddings—LLMs are natively suited to decomposed output and thus were not evaluated under the direct strategy), we compare:
\begin{itemize}
    \item \textbf{Decomposed}: Train a separate classifier for each of WHAT, WHY, HOW; sum predictions for total score.
    \item \textbf{Direct}: Train a single classifier on the total score label.
\end{itemize}
We report accuracy, macro-F1, mean absolute error (MAE), and quadratic weighted kappa (QWK) on the held-out test set ($n = 58$, 80/20 stratified split).

\paragraph{Results.}

\begin{table}[H]
\centering
\caption{Decomposed vs.\ Direct Scoring on Total Score Prediction ($n=58$)}
\label{tab:exp1}
\begin{tabular}{llcccc}
\toprule
\textbf{Approach} & \textbf{Strategy} & \textbf{Accuracy} & \textbf{Macro-F1} & \textbf{MAE} & \textbf{QWK} \\
\midrule
TF-IDF     & Direct     & 0.793 & 0.607 & 0.362 & 0.876 \\
TF-IDF     & Decomposed & 0.810 & 0.540 & 0.259 & 0.954 \\
Embeddings & Direct     & 0.741 & 0.483 & 0.345 & 0.919 \\
Embeddings & Decomposed & 0.810 & 0.497 & 0.224 & 0.968 \\
\bottomrule
\end{tabular}
\end{table}

\paragraph{Discussion.} Decomposed scoring consistently outperforms direct scoring on QWK and MAE for both approaches. The QWK improvement is particularly meaningful: decomposed TF-IDF improves from 0.876 to 0.954, and decomposed embeddings from 0.919 to 0.968. The lower MAE of the decomposed approach indicates that it makes fewer large errors on the total score—precisely the property most important for fair grading. Macro-F1 is slightly lower for decomposed approaches, suggesting that decomposition may reduce coverage of minority classes in the total-score distribution, but the ordinal error reduction takes priority for this task. \textbf{We adopt decomposed scoring for all subsequent analyses and for the final pipeline.}

\subsection{Experiment 2: Cross-Approach and Cross-Model Evaluation}

\paragraph{Motivation.} With decomposed scoring established as superior, we now ask: which scoring approach is best per dimension? We compare all approaches on the same test set to determine which model or technique produces the most accurate per-dimension scores.

\paragraph{Setup.} All approaches are evaluated on the expanded synthetic dataset ($n = 58$). For LLMs, reflections are scored at temperature 0 for deterministic output.

\paragraph{Results.}

\begin{table}[H]
\centering
\caption{Per-Dimension Performance Across Scoring Approaches ($n=58$)}
\label{tab:exp2}
\begin{tabular}{lcccccc}
\toprule
\multirow{2}{*}{\textbf{Model / Approach}} & \multicolumn{2}{c}{\textbf{WHAT}} & \multicolumn{2}{c}{\textbf{WHY}} & \multicolumn{2}{c}{\textbf{HOW}} \\
\cmidrule(lr){2-3} \cmidrule(lr){4-5} \cmidrule(lr){6-7}
 & \textbf{Acc} & \textbf{QWK} & \textbf{Acc} & \textbf{QWK} & \textbf{Acc} & \textbf{QWK} \\
\midrule
TF-IDF + LR        & 0.914 & 0.796 & 0.897 & 0.887 & 0.914 & 0.947 \\
Embeddings + LR    & 0.948 & 0.793 & 0.862 & 0.906 & 0.914 & 0.945 \\
\midrule
Claude Haiku       & \textbf{0.983} & 0.960 & 0.931 & 0.952 & 0.897 & 0.933 \\
Claude Sonnet      & \textbf{0.983} & \textbf{0.960} & 0.862 & 0.907 & 0.914 & 0.943 \\
GPT-4o             & 0.897 & 0.838 & \textbf{0.948} & \textbf{0.963} & 0.862 & 0.913 \\
Llama 3.3-70B      & \textbf{0.983} & \textbf{0.960} & 0.862 & 0.910 & 0.845 & 0.903 \\
Llama 4 Scout      & 0.948 & 0.878 & \textbf{0.966} & \textbf{0.975} & \textbf{0.914} & \textbf{0.944} \\
\bottomrule
\end{tabular}
\end{table}

\paragraph{Discussion.} All LLM-based approaches substantially outperform TF-IDF and embedding baselines on WHAT (QWK 0.796--0.793 vs.\ 0.960 for top LLMs), confirming that surface lexical patterns are insufficient for capturing the specificity of student descriptions. For WHY and HOW, the gap is smaller but LLMs remain consistently stronger. Within the LLM family, Claude Sonnet achieves the highest or tied-highest QWK on WHAT (0.960), remains competitive on WHY (0.907) and HOW (0.943), and provides rich structured evidence strings. Llama 4 Scout achieves the highest WHY QWK (0.975), but at the cost of requiring a separate API and introducing additional operational complexity. These results motivate a Claude Sonnet primary pipeline with Llama 3.3-70B as the cross-model validator (see Section~\ref{sec:pipeline}).

% ─────────────────────────────────────────────
\section{Final Pipeline}
\label{sec:pipeline}
% ─────────────────────────────────────────────

Based on the results of Experiments 1 and 2, our final pipeline uses \textbf{decomposed scoring with Claude Sonnet as the primary grader across all three dimensions (WHAT, WHY, HOW)}. The decomposed strategy was selected based on Experiment~1 results showing consistent QWK and MAE improvements over direct scoring. Claude Sonnet was selected as the single model for all dimensions based on its consistently strong performance across all three rubric dimensions in Experiment~2, its ability to produce rich, structured JSON output with evidence spans, and the practical advantage of using a single provider for the entire scoring pass.

\subsection{Evidence Extraction}

Each score is grounded in a text span extracted directly from the student's reflection. For all three dimensions, the LLM prompt instructs Claude Sonnet to return an exact verbatim text span from the reflection as evidence for the assigned score. A post-hoc verification step checks that each returned evidence string is a strict substring of the original reflection text. Evidence spans that fail verification are flagged for manual review, ensuring auditability and guarding against model hallucination.

\subsection{Cross-Model Consistency Validation}

To provide an independent reliability check on Claude Sonnet's scores, we run \textbf{Llama~3.3-70B} as a secondary grader on the full scored dataset. Llama~3.3-70B was selected for this role for three reasons: (1) it is from a different model family (Meta open-weight vs.\ Anthropic proprietary), providing genuine cross-family independence; (2) it achieved competitive per-dimension performance in Experiment~2 (QWK: WHAT 0.960, WHY 0.910, HOW 0.903); and (3) as an open-weight model, it represents a pathway toward future on-premise deployment for privacy-sensitive settings.

We compute pairwise exact agreement rate and Cohen's weighted kappa per dimension (WHAT, WHY, HOW separately) between Claude Sonnet and Llama~3.3-70B. Reflections where any dimension falls below 0.9 agreement are exported as a CSV containing all predicted scores and evidence from both models, flagged for manual instructor review rather than silently resolved.

\begin{table}[H]
\centering
\caption{Cross-Model Agreement: Claude Sonnet vs.\ Llama~3.3-70B}
\label{tab:consistency}
\begin{tabular}{lcc}
\toprule
\textbf{Dimension} & \textbf{Exact Agreement (\%)} & \textbf{Weighted $\kappa$} \\
\midrule
WHAT & -- & -- \\
WHY  & -- & -- \\
HOW  & -- & -- \\
\bottomrule
\end{tabular}
\end{table}

\noindent\textit{Note: Cross-model consistency results to be filled in upon completion of full pipeline run.}

% ─────────────────────────────────────────────
\section{Discussion}
% ─────────────────────────────────────────────

Our experiments confirm that LLM-based scoring substantially outperforms classical approaches (TF-IDF, sentence embeddings) for rubric-aligned reflection assessment, particularly on the WHAT dimension where the gap in QWK exceeds 0.16 points. This indicates that the specificity and concreteness of student learning descriptions cannot be adequately captured by surface lexical patterns or aggregate semantic similarity—both of which lack the contextual reasoning needed to distinguish specific from generic descriptions.

The decision to use Claude Sonnet for all three dimensions, rather than mixing models, was driven by both performance and practicality. While Llama 4 Scout achieves the highest WHY QWK (0.975 vs.\ Claude Sonnet's 0.907), the dataset is small enough that this gap may not be reliable. More importantly, using a single model simplifies deployment, ensures consistent evidence extraction quality, and avoids the need to maintain multiple API integrations for production use. Claude Sonnet's performance across all three dimensions is consistently strong, and the evidence strings it produces tend to be more precisely grounded in the rubric language.

Llama~3.3-70B was chosen as the cross-model validator specifically because it is from a different model family. Using a second Claude model (e.g., Claude Haiku) for cross-model validation would not reveal cross-provider disagreements that might indicate ambiguous reflections. Llama~3.3-70B's competitive performance (QWK 0.960 on WHAT, 0.910 on WHY, 0.903 on HOW) makes it a meaningful validator—low agreement between Claude Sonnet and Llama~3.3-70B is a genuine signal that a reflection is genuinely ambiguous, not merely an artifact of one model being poorly calibrated for the task.

% ─────────────────────────────────────────────
\section{Limitations}
\label{sec:limitations}
% ─────────────────────────────────────────────

\paragraph{Data.}
We were unable to use a real student reflection dataset due to privacy and IRB constraints. All experimentation was conducted on synthetic datasets. While the expanded synthetic dataset was constructed with faculty guidance and validation to reflect realistic reflection patterns, it was not collected from actual students, and our results may not generalize to authentic writing distributions. Both datasets are small, limiting the statistical power of our evaluation and the reliability of cross-model comparisons.

\paragraph{Synthetic validation gap.}
Because both datasets are synthetic, they may share systematic vocabulary and structural patterns that inflate agreement between all models evaluated on them. A model that overfits to synthetic language patterns may score well here but perform differently on real student text.

\paragraph{Model coverage.}
Our model comparisons were constrained by access and cost. API costs were borne out of pocket, and free-tier quotas limited the number of reflections scored with some models. Cross-model consistency analysis is therefore limited in scope; a more rigorous reliability study would evaluate all models on the full dataset and report confidence intervals around agreement statistics.

% ─────────────────────────────────────────────
\section{Future Work}
% ─────────────────────────────────────────────

\paragraph{Real student data.}
The most impactful next step is collecting and properly anonymizing a corpus of real student reflections with expert-assigned rubric scores. This would enable validation of the pipeline against human raters and meaningful assessment of inter-rater reliability.

\paragraph{Expanded model coverage.}
As API costs decline and open-weight models improve, future work should evaluate a broader set of models and include more robust cross-family agreement analyses with confidence intervals.

\paragraph{Privacy-preserving deployment.}
For deployment in an institutional setting, student data privacy is paramount. Future work should explore on-premise deployment of open-weight models (e.g., Llama via Ollama) eliminating the need to transmit student text to third-party APIs—a direction already motivated by our inclusion of Llama~3.3-70B in the current pipeline.

\paragraph{Instructor-facing application.}
A natural extension is a web application allowing instructors to upload reflection CSVs, review scores and evidence, override predictions, and export annotated results—reducing manual grading burden while preserving instructor oversight.

\paragraph{Fine-tuning on domain data.}
Once sufficient real-data annotations are available, fine-tuning a smaller model on pharmacy reflection data could provide a cost-effective alternative to repeated Claude API calls at scale.

% ─────────────────────────────────────────────
\section{Conclusion}
% ─────────────────────────────────────────────

We presented an automated scoring pipeline for pharmacy student reflections using a decomposed LLM architecture with Claude Sonnet as the primary grader across all three rubric dimensions (WHAT, WHY, HOW). Through systematic experimentation across TF-IDF, embedding, and LLM approaches, we demonstrated that LLM-based scoring substantially outperforms classical methods, particularly for capturing the specificity of student descriptions. Experiment~1 established that decomposed scoring consistently outperforms direct total-score prediction, and Experiment~2 identified Claude Sonnet as the strongest single-model choice for deployment. Evidence extraction grounds each prediction in the student's own words, and cross-model validation with Llama~3.3-70B surfaces ambiguous cases for instructor review.

The primary limitation of this work is data: privacy constraints prevented use of real student reflections, and our synthetic datasets are small. Despite these constraints, the system demonstrates a viable, interpretable architecture for rubric-aligned reflection scoring that is ready for validation once real-data annotation becomes feasible.

% ─────────────────────────────────────────────
\bibliographystyle{plainnat}
\bibliography{references}

% ─────────────────────────────────────────────
\appendix
\section{Course Rubric for Reflection Assessment}
\label{app:rubric}
% ─────────────────────────────────────────────

The following rubric was provided by faculty at the University of Washington School of Pharmacy and is reproduced here with their approval. This rubric forms the basis of the WHAT, WHY, and HOW scoring dimensions used throughout this work.

\vspace{1em}

\begin{table}[H]
\centering
\caption{UW School of Pharmacy Reflection Rubric (Faculty-Validated)}
\label{tab:rubric}
\begin{tabular}{p{2.2cm}p{3.8cm}p{3.8cm}p{3.8cm}}
\toprule
\textbf{Dimension} & \textbf{Score 2} & \textbf{Score 1} & \textbf{Score 0} \\
\midrule
\textbf{WHAT}\newline(Description) &
Clear, specific description of a concrete learning experience, skill, or concept with relevant supporting details. &
Description is present but vague, incomplete, or lacking sufficient specificity. &
No meaningful description of learning, or content is off-topic. \\
\midrule
\textbf{WHY}\newline(Analysis) &
Clear explanation of why the learning is personally significant, with insight, self-awareness, and concrete examples. &
Analysis is surface-level, generic, or lacks a personal connection to the experience. &
No meaningful analysis of why the learning matters. \\
\midrule
\textbf{HOW}\newline(Integration) &
Specific, actionable, and realistic plan for applying the learning in future pharmacy practice. &
Plan is vague, generic, or not clearly actionable in a pharmacy context. &
No mention of future application, or application is purely hypothetical. \\
\bottomrule
\end{tabular}
\end{table}

\noindent\textit{Note: The rubric above reflects the scoring criteria as validated by faculty collaborators. Exact rubric language is reproduced with professor approval and may be updated if the course rubric is revised in future iterations.}

\end{document}
